% Ampoule à décanter : extraction du diiode par le cyclohexane (exercice 14, chapitre 7)
\begin{tikzpicture}[x=1cm,y=1cm]
\def\ampoule{(-0.18,3.6) -- (-0.18,3.3) .. controls (-0.18,3.0) and (-1.05,2.7) .. (-1.05,1.75)
  .. controls (-1.05,1.0) and (-0.35,0.62) .. (-0.12,0.42) -- (-0.12,0.1) -- (0.12,0.1) -- (0.12,0.42)
  .. controls (0.35,0.62) and (1.05,1.0) .. (1.05,1.75) .. controls (1.05,2.7) and (0.18,3.0) .. (0.18,3.3) -- (0.18,3.6)}
\begin{scope}
\clip \ampoule -- cycle;
\fill[pcJaune] (-1.2,0) rectangle (1.2,1.55);
\fill[magenta!30!violet!35!white] (-1.2,1.55) rectangle (1.2,2.45);
\end{scope}
\draw[pcSurface] (-1.04,1.55) -- (1.04,1.55);
\draw[pcVerre] \ampoule;
\filldraw[pcAppareil] (-0.24,3.6) rectangle (0.24,3.95);
\filldraw[pcAppareil] (-0.32,-0.02) rectangle (0.32,0.16);
\draw[otNavy,line width=0.6pt] (-0.06,-0.02) -- (-0.06,-0.45) (0.06,-0.02) -- (0.06,-0.45);
\pcNote[west]{(0.85,2.0)}{(1.6,2.35)}{phase organique : cyclohexane (densité 0,78), diiode extrait}
\pcNote[west]{(0.9,1.1)}{(1.6,1.0)}{phase aqueuse (densité 1,0)}
\pcNote[west]{(0.32,0.07)}{(1.6,0.2)}{robinet}
\end{tikzpicture}
\par\smallskip\legende{Ampoule à décanter après agitation puis décantation : le diiode extrait colore la phase organique en violet.}
